\chapter{Goals}

The goal of the project is to develop the ground and on board side of OSI Layer 2 of ERWIN’s TM chain. As already specified in the scope, the goal is to be compliant with CCSDS 132.0-B-3 and CCSDS 131.0-B-5.

\section{CCSDS 132.0-B-3}

The current design aims to implement full first header support. The secondary header is not going to be implemented. Likewise, neither the OCF nor the FECF parts of the trailer are going to be implemented. Insertion of idle data, to keep a consistent data rate, is part of the current goal. Only one master channel, but multiple virtual channels will be supported, if support for multiple channels is required by ERWIN’s system design.

\section{CCSDS 131.0-B-5}

As the FECF is not going to be implemented the following measures for error correction are going  to be developed in the lower layer: Reed Solomon will form the inner error correction method, with the following parameters: $E=16$, $I=1$ and virtual fill disabled. Furthermore, the outer error correction method will consist of a convolutional code of code rate $r = 50\%$. Between these layers the suggested randomization, with polynomial of degree 17 and the required ASM field will also be implemented.

\section{Test System}

Additionally a full hardware in the loop test system is planned. For that purpose 3 additional systems are going to be implemented.

\subsection{OBC Emulator}

A stub On-Board-Computer will be developed, to allow sending fake space packets from the processing system (PS) of the target SoC to the CCSDS 132.0-B-3 layer, implemented on the PL. 

\subsection{AXI to ETH Interface}

To allow routing fully processed TM frames out of the OnBoard SoC, an interface will be developed, to send out the frames via an Ethernet interface, in the same format, as the ground system requires them.

\subsection{Control Software}

Finally a control software will be developed, which will manage generating the fake space packets for the OBC emulator and also receive the final unpacked space packets decoded by the ground system. Therefore, the control software will be able to monitor packet loss, latency (to some degree) and throughput of the final system. Additionally the software will be able to intercept packets generated by the AXI to ETH interface, to simulate upsets in the transmission to also test the error correction’s performance.

\section{Additional Goals}

If all previously specified goals are met, additional goals will be implemented, to accommodate a broader usage of the developed product. These additional features will focus on more CCSDS 131.0-B-5 features, to allow a mission specific configuration of data flow. Specifically these goals are the integration of a secondary header into the TM frame, as well as a cyclic redundancy check (CRC) checksum in the FECF. Lastly the OCF will be made available, but will not be filled with meaningful data, as this is up to the mission specification.